\section{时空基本假设}
\subsection{狭义相对论}
一些约定
\begin{enumerate}
    \item  采用高能物理中标准的自然单位制：
      \begin{equation}\label{eq:natural-units}
        \hbar = c = 1
      \end{equation}
      在此单位制下，所有物理量的量纲都可以用质量（或等价地，能量、逆长度、逆时间）的幂次表达：
      \[
        [E] = [m] = [p] = [1/L] = [1/T] = \text{mass}
      \]
      常用的量纲关系为：
      \begin{center}
      \begin{tabular}{lll}
        截面 & $[\sigma] = [\text{mass}]^{-2}$ & $1\;\text{GeV}^{-2} \approx 0.389\;\text{mb}$ \\
        寿命 & $[\tau] = [\text{mass}]^{-1}$ & $1\;\text{GeV}^{-1} \approx 6.58\times10^{-25}\;\text{s}$ \\
        长度 & $[L] = [\text{mass}]^{-1}$   & $1\;\text{GeV}^{-1} \approx 0.197\;\text{fm}$
      \end{tabular}
      \end{center}
      需要恢复 $\hbar$ 和 $c$ 时，可利用量纲分析唯一确定因子。
    \item  度规约定
      \begin{equation}\label{eq:metric}
        \eta_{\mu\nu} = \eta^{\mu\nu}
        = \mathrm{diag}(+1,\,-1,\,-1,\,-1)
      \end{equation}
    \item 
        闵氏时空中两个事件之间的\textbf{不变间隔}为：
          \begin{equation}\label{eq:interval}
            ds^2 = \eta_{\mu\nu}\,dx^\mu\,dx^\nu
                = (dx^0)^2 - (dx^1)^2 - (dx^2)^2 - (dx^3)^2
                = dt^2 - d\mathbf{x}^2
          \end{equation}
          间隔的分类：
          \begin{itemize}
            \item $ds^2 > 0$：\textbf{类时}（timelike）——
                  两个事件可以被同一观测者依次经历
            \item $ds^2 < 0$：\textbf{类空}（spacelike）——
                  两个事件不能有因果联系
            \item $ds^2 = 0$：\textbf{类光}（lightlike / null）——
                  光信号恰好连接两个事件
          \end{itemize}
    \item 指标约定：
          \begin{itemize}
            \item 希腊字母 $\mu, \nu, \rho, \sigma, \ldots = 0, 1, 2, 3$ 遍历时空全部四个分量。
            \item 拉丁字母 $i, j, k, \ldots = 1, 2, 3$ 仅遍历空间三个分量。
            \item 对重复的上下指标自动求和（\textbf{Einstein 求和约定}）。
          \end{itemize}

          \noindent\textbf{指标升降：} 度规 $\eta_{\mu\nu}$ 及其逆 $\eta^{\mu\nu}$
          （在当前约定下二者数值相同）用于升降指标：
          \begin{equation}
            A_\mu = \eta_{\mu\nu}A^\nu\,,\qquad A^\mu = \eta^{\mu\nu}A_\nu
          \end{equation}
          对于逆变（上指标）四矢量 $A^\mu = (A^0, \mathbf{A})$，
          对应的协变（下指标）四矢量为 $A_\mu = (A^0, -\mathbf{A})$。
    \item 常用不变量：
      \begin{itemize}
    \item 质壳条件：$p^\mu p_\mu = p^2 = E^2 - |\mathbf{p}|^2 = m^2$
    \item 四维散度（d'Alembert 算符）：
          $\partial^\mu\partial_\mu
          = \partial_t^2 - \nabla^2
          \equiv \Box$
    \item 四维体积元：$d^4x = dt\,d^3\mathbf{x}$
          （Lorentz 标量\footnote{
          严格地说，$d^4x$ 在固有正时 Lorentz 变换下不变
          （因为 $|\det\Lambda| = 1$），
          但在包含空间或时间反演的变换下可能获得符号。}）
      \end{itemize}
    \item 完全反对称张量:
    定义四维 Levi-Civita 符号\footnote{
    注意区分 Levi-Civita \textbf{符号}（不是张量）$\epsilon^{\mu\nu\rho\sigma}$
    和 Levi-Civita \textbf{张量}
    $\varepsilon^{\mu\nu\rho\sigma} = \epsilon^{\mu\nu\rho\sigma}/\sqrt{|g|}$。
      在闵氏时空中 $\sqrt{|g|} = 1$，二者数值相同。
    }：
    \begin{equation}
  \epsilon^{0123} = +1\,,\qquad
  \epsilon^{\mu\nu\rho\sigma}\text{完全反对称}
    \end{equation}
    降指标后：$\epsilon_{0123} = \eta_{0\alpha}\eta_{1\beta}\eta_{2\gamma}\eta_{3\delta}
      \,\epsilon^{\alpha\beta\gamma\delta} = -1$。
常用恒等式：
    \begin{equation}
  \epsilon^{\mu\nu\rho\sigma}\epsilon_{\mu\nu\rho\sigma} = -24\,,\qquad
  \epsilon^{\mu\nu\rho\sigma}\epsilon_{\mu\nu\rho\tau} = -6\,\delta^\sigma{}_\tau
    \end{equation}
\end{enumerate}
\subsection{Lorentz 变换与 Poincar\'e 变换}

本节系统地建立 Lorentz 变换和 Poincar\'e 变换的数学结构。这些变换构成的群是
相对论量子场论的\textbf{时空对称群}
\topictitle{Lorentz 变换}
\textbf{Lorentz 变换}是闵氏时空中保持不变间隔 $ds^2$ 不变的线性变换：
\begin{equation}\label{eq:Lorentz-def}
  x'^\mu = \Lambda^\mu{}_\nu\, x^\nu
\end{equation}
要求 $ds'^2 = ds^2$，即 $\eta_{\mu\nu}\,dx'^\mu\,dx'^\nu = \eta_{\rho\sigma}\,dx^\rho\,dx^\sigma$，
代入 $dx'^\mu = \Lambda^\mu{}_\rho\,dx^\rho$，由 $dx^\rho$ 的任意性得到保度规条件：
\begin{equation}\label{eq:metric-preserve}
  \eta_{\mu\nu}\,\Lambda^\mu{}_\rho\,\Lambda^\nu{}_\sigma = \eta_{\rho\sigma}
\end{equation}或用矩阵记号：$\Lambda^T \eta\,\Lambda = \eta$。
\vspace{2ex}\hrule
\vspace{2ex}
从保度规条件 (\ref{eq:metric-preserve}) 可以推导出两个重要的约束：\begin{itemize}
  \item \noindent\textbf{1. 行列式约束。}
          对 (\ref{eq:metric-preserve}) 两端取行列式：
          \begin{equation}
            \det(\Lambda^T\eta\Lambda) = \det\eta
            \quad\Longrightarrow\quad
            (\det\Lambda)^2 \cdot \det\eta = \det\eta
            \quad\Longrightarrow\quad
            (\det\Lambda)^2 = 1
          \end{equation}
          因此
          \begin{equation}\label{eq:det-constraint}
            \det\Lambda = +1\quad\text{（固有变换）}
            \qquad\text{或}\qquad
            \det\Lambda = -1\quad\text{（非固有变换）}
          \end{equation}
  \item \noindent\textbf{2. 时间分量约束。}
        取 (\ref{eq:metric-preserve}) 中 $\rho = \sigma = 0$：
        \begin{equation}
          \eta_{\mu\nu}\Lambda^\mu{}_0\Lambda^\nu{}_0 = \eta_{00} = 1
          \quad\Longrightarrow\quad
          (\Lambda^0{}_0)^2 - \sum_{i=1}^{3}(\Lambda^i{}_0)^2 = 1
        \end{equation}
        由于 $\sum_i(\Lambda^i{}_0)^2 \geqslant 0$，必然有
        \begin{equation}\label{eq:time-constraint}
          (\Lambda^0{}_0)^2 \geqslant 1
          \quad\Longrightarrow\quad
          \Lambda^0{}_0 \geqslant +1\quad\text{（正时变换）}
          \qquad\text{或}\qquad
          \Lambda^0{}_0 \leqslant -1\quad\text{（非正时变换）}
        \end{equation}
  \item 行列式和 $\Lambda^0{}_0$ 的符号都是离散的，因此它们在
            Lorentz 群的连续变换下不可能改变（不可能从 $+1$ 连续变到 $-1$）。
            这将 Lorentz 群 $O(1,3)$ 分为四个互不连通的分支：

            \begin{center}
            \renewcommand{\arraystretch}{1.5}
            \begin{tabular}{ccccl}
              \hline
              分支 & $\det\Lambda$ & $\Lambda^0{}_0$ & 名称 & 典型元素 \\
              \hline
              $\mathcal{L}^\uparrow_+$ & $+1$ & $\geqslant+1$
                & 固有正时 & 恒等 $\mathbf{1}$，旋转 $R$，boost $L$ \\
              $\mathcal{L}^\uparrow_-$ & $-1$ & $\geqslant+1$
                & 非固有正时 & 空间反演 $P$ \\
              $\mathcal{L}^\downarrow_-$ & $-1$ & $\leqslant-1$
                & 非固有非正时 & 时间反演 $T$ \\
              $\mathcal{L}^\downarrow_+$ & $+1$ & $\leqslant-1$
                & 固有非正时 & 时空反演 $PT$ \\
              \hline
            \end{tabular}
            \end{center}
            \noindent 其中离散变换的矩阵形式为：
            \begin{equation}\label{eq:P-T-matrices}
              P^\mu{}_\nu = \mathrm{diag}(+1,-1,-1,-1)\,,\qquad
              T^\mu{}_\nu = \mathrm{diag}(-1,+1,+1,+1)
            \end{equation}
            四个分支的关系为：
            \[
              \mathcal{L}^\uparrow_- = P\cdot\mathcal{L}^\uparrow_+\,,\qquad
              \mathcal{L}^\downarrow_- = T\cdot\mathcal{L}^\uparrow_+\,,\qquad
              \mathcal{L}^\downarrow_+ = PT\cdot\mathcal{L}^\uparrow_+
            \]
            只有\textbf{固有正时 Lorentz 群} $\mathcal{L}^\uparrow_+ = SO^+(1,3)$
            连续包含恒等变换，因此它是 $O(1,3)$ 的\textbf{单位连通分支}。
            在量子场论中，连续对称性由 $SO^+(1,3)$ 描述；离散对称性 $P$、$T$
            及其组合需要单独讨论\footnote{$C$（电荷共轭）是内部对称性，不属于时空变换群。
            CPT 定理则断言：任何满足 Lorentz 不变性、局域性和幺正性的
            量子场论都自动具有 $CPT$ 对称性。}
\end{itemize}
\vspace{3ex}

\begin{center}
  \large 有限 Lorentz 变换的显式形式
\end{center}
\vspace{2ex}
$SO^+(1,3)$ 的一般元素可以参数化为\textbf{旋转}和\textbf{boost}的组合。
以下给出两类基本变换的矩阵形式。
\begin{itemize}
  \item \textbf{空间旋转。}
    以绕 $x^3$（$z$）轴旋转角度 $\theta$ 为例：
    \begin{equation}\label{eq:rotation-z}
      R_3(\theta) =
      \begin{pmatrix}
        1 & 0 & 0 & 0\\
        0 & \cos\theta & -\sin\theta & 0\\
        0 & \sin\theta & \cos\theta & 0\\
        0 & 0 & 0 & 1
      \end{pmatrix}
    \end{equation}
    三个空间旋转 $R_1(\theta_1)$、$R_2(\theta_2)$、$R_3(\theta_3)$
    生成三维旋转群 $SO(3) \subset SO^+(1,3)$，有 3 个参数。
  \item \textbf{Lorentz boost。}
以沿 $x^3$（$z$）方向的 boost 为例，引入快度（rapidity）$\phi$\footnote{
快度的物理含义：
$\tanh\phi = v$（$c=1$），
$\cosh\phi = \gamma = 1/\sqrt{1-v^2}$，
$\sinh\phi = \gamma v$。
快度的最大优势在于同向 boost 的快度可直接相加：
$L(\phi_2)\,L(\phi_1) = L(\phi_1 + \phi_2)$，
而速度的相加则需要使用相对论速度合成公式。}：
\begin{equation}\label{eq:boost-z}
  L_3(\phi) =
  \begin{pmatrix}
    \cosh\phi & 0 & 0 & \sinh\phi\\
    0 & 1 & 0 & 0\\
    0 & 0 & 1 & 0\\
    \sinh\phi & 0 & 0 & \cosh\phi
  \end{pmatrix}
\end{equation}
注意 boost 矩阵\textbf{不是正交的}（$L^T L \neq \mathbf{1}$，
而是 $L^T\eta L = \eta$），
这正是 5.2 中讨论的 Lorentz 群非紧致性的具体体现。

沿任意方向 $\hat{\mathbf{n}}$ 的一般 boost 为\footnote{
此公式在 5.3.5 中计算矢量场极化矢量时有直接应用。}：
\begin{equation}\label{eq:general-boost}
  \Lambda^0{}_0 = \cosh\phi\,,\quad
  \Lambda^0{}_i = \Lambda^i{}_0 = \hat{n}_i\sinh\phi\,,\quad
  \Lambda^i{}_j = \delta_{ij} + (\cosh\phi - 1)\hat{n}_i\hat{n}_j
\end{equation}
三个独立的 boost 方向提供 3 个参数。
连同 3 个旋转参数，$SO^+(1,3)$ 共有 \textbf{6 个参数}。
\end{itemize}
\topictitle{Poincar\'e 变换}
闵氏时空的完整对称群不仅包括 Lorentz 变换（保持原点不动），还包括时空平移。
二者的组合称为\textbf{Poincar\'e 变换}（也称为非齐次 Lorentz 变换）：
\begin{equation}\label{eq:Poincare-def}
  x'^\mu = \Lambda^\mu{}_\nu\,x^\nu + a^\mu
\end{equation}
其中 $\Lambda \in SO^+(1,3)$，$a^\mu \in \mathbb{R}^{1,3}$。
我们用二元组 $(\Lambda, a)$ 标记一个 Poincar\'e 变换。


\begin{tcolorbox}[colback=white!5,colframe=black!45]
  \textbf{群乘法。}
连续执行两次变换 $(\Lambda_1, a_1)$ 和 $(\Lambda_2, a_2)$：
\begin{align}
  x'' &= \Lambda_2\,x' + a_2 = \Lambda_2(\Lambda_1 x + a_1) + a_2
       = (\Lambda_2\Lambda_1)\,x + (\Lambda_2 a_1 + a_2)
\end{align}
因此群乘法为：
\begin{equation}\label{eq:Poincare-mult}
\boxed{
  (\Lambda_2, a_2)\cdot(\Lambda_1, a_1)
  = (\Lambda_2\Lambda_1,\;\Lambda_2 a_1 + a_2)
}
\end{equation}
验证群公理：
\begin{itemize}
  \item \textbf{封闭性}：由 (\ref{eq:Poincare-mult}) 显然。
  \item \textbf{结合律}：矩阵乘法和向量加法的结合律保证。
  \item \textbf{单位元}：$(\mathbf{1}, 0)$。
  \item \textbf{逆元}：$(\Lambda, a)^{-1} = (\Lambda^{-1},\,-\Lambda^{-1}a)$\footnote{
        验证：$(\Lambda^{-1}, -\Lambda^{-1}a)\cdot(\Lambda, a)
        = (\Lambda^{-1}\Lambda,\;\Lambda^{-1}a - \Lambda^{-1}a) = (\mathbf{1}, 0)$。
        这正是 \S3.2 中推导 $U(\Lambda,a)U(1+\omega,\epsilon)U^\dagger(\Lambda,a)$
        时所使用的逆元形式。}。
\end{itemize}
\end{tcolorbox}
\vspace{2ex}
\par
Poincar\'e 群记为 $\mathrm{ISO}^+(1,3)$，它具有半直积结构：
\begin{equation}\label{eq:semidirect}
\boxed{
  \mathrm{ISO}^+(1,3) = \mathbb{R}^{1,3} \rtimes SO^+(1,3)
}
\end{equation}
这意味着：
\begin{enumerate}
  \item 平移子群 $\mathbb{R}^{1,3}$（对应 $(\mathbf{1}, a)$）是\textbf{正规子群}
        （不变子群），因为
        $(\Lambda, 0)\,(\mathbf{1}, a)\,(\Lambda^{-1}, 0) = (\mathbf{1}, \Lambda a) \in \mathbb{R}^{1,3}$。
  \item Lorentz 子群 $SO^+(1,3)$（对应 $(\Lambda, 0)$）\textbf{不是}正规子群。
  \item 群乘法不是简单的分量相乘（$a_2$ 被 $\Lambda_2$ 作用），
        所以不是直积而是半直积。
\end{enumerate}
Poincar\'e 群有 $6 + 4 = \mathbf{10}$ 个参数：
6 个 Lorentz 参数（3 旋转 + 3 boost）加上 4 个平移参数。
相应地，Poincar\'e 代数有 10 个生成元。
\begin{center}
  \large ***
\end{center}
在 1.5.1 中我们讨论了投影表示和万有覆盖群的概念。
三维旋转群 $SO(3)$ 的万有覆盖群是 $SU(2)$，
二者之间存在 2:1 的覆盖映射。完全平行地：

\begin{center}
\fbox{\parbox{0.8\textwidth}{
固有正时 Lorentz 群 $SO^+(1,3)$ 的万有覆盖群是
$SL(2,\mathbb{C})$——所有行列式为 1 的 $2\times 2$ 复矩阵构成的群。
覆盖映射是 2:1 的：$SL(2,\mathbb{C})$ 中的 $A$ 和 $-A$
对应 $SO^+(1,3)$ 中的同一个 Lorentz 变换。
}}
\end{center}

\noindent 覆盖映射的显式构造如下\footnote{
这一构造将在 5.2 中从 Lorentz 代数的
$\mathfrak{su}(2)\oplus\mathfrak{su}(2)$ 分解中自然出现。}。
定义从四维坐标到 $2\times 2$ 厄米矩阵的映射：
\begin{equation}\label{eq:x-to-matrix}
  \bar{x} \;\equiv\; x^\mu\sigma_\mu
  = x^0\,\mathbf{1} + x^i\sigma_i
  = \begin{pmatrix}
      x^0 + x^3 & x^1 - ix^2 \\
      x^1 + ix^2 & x^0 - x^3
    \end{pmatrix}
\end{equation}
其中 $\sigma_\mu = (\mathbf{1}, \sigma^1, \sigma^2, \sigma^3)$，$\sigma^i$ 是 Pauli 矩阵。
关键性质是 $\det\bar{x} = (x^0)^2 - |\mathbf{x}|^2 = x^\mu x_\mu$，即矩阵的行列式恰好
等于 Lorentz 不变量。

对于 $A \in SL(2,\mathbb{C})$，定义变换：
\begin{equation}\label{eq:SL2C-action}
  \bar{x}' = A\,\bar{x}\,A^\dagger
\end{equation}
由于 $\det\bar{x}' = |\det A|^2\det\bar{x} = \det\bar{x}$
（因为 $\det A = 1$），
且 $\bar{x}'$ 仍然是厄米矩阵（因此可以展开为 $x'^\mu\sigma_\mu$），
这定义了一个保持 Lorentz 不变量的线性变换 $x^\mu \to x'^\mu$，即一个 Lorentz 变换
$\Lambda(A)$。容易验证 $\Lambda(A) = \Lambda(-A)$，因此映射 $A \mapsto \Lambda(A)$ 恰好是 2:1 的：
\begin{equation}
  SL(2,\mathbb{C})\big/\mathbb{Z}_2 \;\cong\; SO^+(1,3)
\end{equation}
由于 $SL(2,\mathbb{C})$ 是单连通的\footnote{
$SL(2,\mathbb{C})$ 作为实流形同胚于 $S^3 \times \mathbb{R}^3$，是单连通的。
而 $SO^+(1,3)$ 是双连通的（$\pi_1(SO^+(1,3)) = \mathbb{Z}_2$），
其基本群恰好对应覆盖映射的核 $\{+\mathbf{1}, -\mathbf{1}\}$。
}，$SL(2,\mathbb{C})$ 就是 $SO^+(1,3)$ 的万有覆盖群。
这一结构的物理后果是深远的：
\begin{itemize}
  \item 为了描述自然界中所有可能的粒子，必须使用 $SL(2,\mathbb{C})$ 的表示，
        而非 $SO^+(1,3)$ 自身的表示——这正是\textbf{旋量表示}
        （5.2 中的 $(A,B)$ 表示）存在的数学根源。
  \item 半整数自旋粒子（费米子）对应于 $SL(2,\mathbb{C})$ 的忠实表示，
        这些表示在 $SO^+(1,3)$ 上是双值的（旋转 $2\pi$ 后态矢量获得负号）。
        这与 1.5.1 中 $SU(2)$ 覆盖 $SO(3)$ 导致半整数自旋的逻辑完全一致。
\end{itemize}